% Paper 3 — Discussion (borrador 2026-10-17).

\section{Discussion}\label{sec:discussion}

\paragraph{What a gate should report.}
A statistical gate makes two choices: the form of its test and the level of its
threshold. The concurrent gates have converged on the first: a paired test is now common to PACE, HarnessBank and SAGE, discordance-based in PACE and SAGE and a $z$-test on per-task differences in HarnessBank \citep{zayxshawn2026pace,luo2026harnessbank,wang2026sage}. Our results say the
second choice matters at least as much: the same rule moves across two orders of
magnitude of false acceptances when only its margin changes, and a central
estimate of the noise leaves about 18\% of null changes committed regardless of
how it is derived. We therefore suggest that a gate report three quantities
alongside its gains: the noise estimate it used and how it was obtained, the
threshold that estimate produced, and the false-acceptance rate of that
threshold on null comparisons of its own configurations. The last one costs
nothing beyond the repeated runs that a noise estimate already requires, since
every pair of repetitions is a null comparison.

\paragraph{Re-derive the threshold for each incumbent.}
Where the noise of configurations differs, a tolerance measured once on the base
is a bet on configurations that have not been measured. Our data show the bet
losing in both directions: it accepts null changes from incumbents noisier than
the base and rejects real improvements from incumbents quieter than it. The
remedy is to re-estimate the noise for the incumbent being challenged, from its
own repetitions, with a tail estimator. That is what SAGE proposes as future
work, ``an adaptive threshold that tightens or loosens $\alpha$ from an online
estimate of task noise'' \citep{wang2026sage}, and our measurements are evidence
that the extension is worth its cost. The cost is runs, not power. Since an
incumbent is evaluated many times over an evolution, its repetitions can be
accumulated across rounds rather than paid at once, so the online version that
SAGE describes is the natural implementation.

\paragraph{When a fixed threshold is enough.}
A single calibration is adequate in two situations. One is when every
configuration of the evolution is about as noisy as the base, for example with
models that are nearly deterministic at the temperature used. The other is when
the base happens to be the noisiest configuration the evolution visits, in which
case the fixed threshold is conservative and pays only in power. Neither can be
known without measuring the incumbents, which is the measurement that a fixed
threshold is meant to avoid. A cheap check is to repeat a few incumbents and
compare their noise with the base's; if they differ, the fixed threshold is
already miscalibrated.

\paragraph{Noise does not vanish with scale in this range.}
Our archive comes from small local models, which a single workstation can serve; PACE also evaluates small models, as a minimal testbed for the accept decision \citep{zayxshawn2026pace}. One might expect the problem to fade as models
grow. In our data the typical configuration does get quieter, but the tail does
not: the median flip rate is 6.3\% for arms of models up to 4B parameters (105
arms), 7.0\% for 7--12B (25) and 2.1\% for 20B and above (17), while the 90th
percentile is 17\%, 21\% and 25\% respectively. With only 17 arms at 20B and above and no test across bands, this is descriptive; if it holds, larger models make a quiet base more likely without making a noisy incumbent less likely, which is the configuration in which a single calibration fails worst (Table~\ref{tab:fixed}). Two arms whose size we could not assign are left out. This complements the findings that the gains of harness evolution are unstable \citep{wu2026harnessdev} and generalize poorly beyond the tasks they were evolved on \citep{wu2026harnessdev,wang2026rethinking}:
part of that instability can come from the accept step, where some of the
committed changes were never improvements.

\paragraph{Future work.}
Three extensions follow. First, a sequential version of the per-configuration
rule, which accumulates an incumbent's repetitions across rounds and combines
them with an anytime-valid test such as PACE's, would remove the up-front cost.
Second, the frontier regime needs the same measurement: whether the noise of
configurations is heterogeneous when the models are large and the temperature is
low is an empirical question that our archive cannot answer. Third, the power
analysis should be repeated with the changes that real optimizers propose, which
alter the noise as well as the mean.
